% GENERATED FILE. Edit canonical lessons, then run npm run generate:books.
% canonical-source-hash: fnv1a64:978b6027b9965be5
% canonical-lessons: SA-C479-vyadhih, SA-C479-madhumehah, SA-C479-ksayah, SA-C479-sancarati, SA-C479-mari, SA-R479-review-472-479

\chapter{Diseases}
\label{ch:sa-diseases}
\hlchaptermodality{479}

\begin{chapteropening}
\textbf{By the end of this chapter:} \emph{I can talk about illnesses that last or spread.}

Complete the last lesson of chapter 479: I can talk about illnesses that last or spread.
\end{chapteropening}

\section[\texorpdfstring{\sk{व्याधिः} (vyādhiḥ)}{व्याधिः (vyādhiḥ)}]{\sk{व्याधिः} (vyādhiḥ) — a disease}

\label{lesson:SA-C479-vyadhih}



\begin{quote}
Before the new one: say the Sanskrit for suffering, then the Sanskrit for unbearable.
\end{quote}

\subsection*{You'll want to know: \sk{व्याधिः}}
\textbf{\sk{व्याधिः}} — \emph{vyādhiḥ} — \textquotedblleft{}a disease\textquotedblright{}.

\begin{quote}
\sk{वृद्धस्य} \sk{व्याधिः} \sk{चिरं} \sk{तिष्ठति।}
\end{quote}

\begin{quote}
\emph{vṛddhasya vyādhiḥ ciraṃ tiṣṭhati.}
\end{quote}

\begin{quote}
The old man's disease lingers.
\end{quote}

A masculine noun.

\begin{cousinweb}
\textbf{\sk{व्याधिः}} (\emph{vyādhiḥ}) is built from \emph{vi}, \emph{ā} and the root √dhā, \textquotedblleft{}to put\textquotedblright{}; MW thinks a link with √vyadh, \textquotedblleft{}to pierce\textquotedblright{}, less likely, and gives \textquotedblleft{}disorder, disease, ailment\textquotedblright{} (MW).
\end{cousinweb}

\begin{grammarlens}[title={What you've built}]
One more word for describing what you lived and felt.
\end{grammarlens}

\subsection*{Your turn}
\begin{itemize}
\raggedright
  \item \emph{Say it:} \emph{vyādhiḥ}
  \item \emph{Say it:} \emph{vyādhiḥ}, once more
  \item \emph{Say it:} \emph{vṛddhasya vyādhiḥ ciraṃ tiṣṭhati.} — The old man's disease lingers.
  \item \emph{Recall:} say the Sanskrit for suffering, then the Sanskrit for unbearable, then say \emph{vyādhiḥ} again
\end{itemize}

\begin{tcolorbox}[breakable,colback=teal!4,colframe=teal!35!black,title={Before you move on}]
What does \sk{व्याधिः} mean? (\textquotedblleft{}A disease\textquotedblright{}.) Say it once more.
\end{tcolorbox}
\section[\texorpdfstring{\sk{मधुमेहः} (madhumehaḥ)}{मधुमेहः (madhumehaḥ)}]{\sk{मधुमेहः} (madhumehaḥ) — diabetes}

\label{lesson:SA-C479-madhumehah}



\begin{quote}
Before the new one: say the Sanskrit for unbearable, then the Sanskrit for a disease.
\end{quote}

\subsection*{You'll want to know: \sk{मधुमेहः}}
\textbf{\sk{मधुमेहः}} — \emph{madhumehaḥ} — \textquotedblleft{}diabetes\textquotedblright{}.

\begin{quote}
\sk{मधुमेहे} \sk{शर्करां} \sk{न} \sk{खादेत्।}
\end{quote}

\begin{quote}
\emph{madhumehe śarkarāṃ na khādet.}
\end{quote}

\begin{quote}
With diabetes one should not eat sugar.
\end{quote}

A masculine noun.

\begin{cousinweb}
\textbf{\sk{मधुमेहः}} (\emph{madhumehaḥ}) is built from \textbf{\sk{मधु}}, \textquotedblleft{}honey\textquotedblright{}, and \emph{meha}, \textquotedblleft{}urine\textquotedblright{}, from the root √mih; MW gives \textquotedblleft{}honey-like urine, diabetes\textquotedblright{} (MW).
\end{cousinweb}

\begin{grammarlens}[title={What you've built}]
One more word for describing what you lived and felt.
\end{grammarlens}

\subsection*{Your turn}
\begin{itemize}
\raggedright
  \item \emph{Say it:} \emph{madhumehaḥ}
  \item \emph{Say it:} \emph{madhumehaḥ}, once more
  \item \emph{Say it:} \emph{madhumehe śarkarāṃ na khādet.} — With diabetes one should not eat sugar.
  \item \emph{Recall:} say the Sanskrit for unbearable, then the Sanskrit for a disease, then say \emph{madhumehaḥ} again
\end{itemize}

\begin{tcolorbox}[breakable,colback=teal!4,colframe=teal!35!black,title={Before you move on}]
What does \sk{मधुमेहः} mean? (\textquotedblleft{}Diabetes\textquotedblright{}.) Say it once more.
\end{tcolorbox}
\section[\texorpdfstring{\sk{क्षयः} (kṣayaḥ)}{क्षयः (kṣayaḥ)}]{\sk{क्षयः} (kṣayaḥ) — a wasting disease}

\label{lesson:SA-C479-ksayah}



\begin{quote}
Before the new one: say the Sanskrit for a disease, then the Sanskrit for diabetes.
\end{quote}

\subsection*{You'll want to know: \sk{क्षयः}}
\textbf{\sk{क्षयः}} — \emph{kṣayaḥ} — \textquotedblleft{}a wasting disease\textquotedblright{}.

\begin{quote}
\sk{क्षयः} \sk{फुप्फुसस्य} \sk{रोगः} \sk{अस्ति।}
\end{quote}

\begin{quote}
\emph{kṣayaḥ phupphusasya rogaḥ asti.}
\end{quote}

\begin{quote}
Consumption is a disease of the lungs.
\end{quote}

A masculine noun.

\begin{cousinweb}
\textbf{\sk{क्षयः}} (\emph{kṣayaḥ}) comes from the root √kṣi, \textquotedblleft{}to waste away\textquotedblright{}; among MW's senses are \textquotedblleft{}loss, waste\textquotedblright{} and \textquotedblleft{}consumption, phthisis\textquotedblright{} (MW).
\end{cousinweb}

\begin{grammarlens}[title={What you've built}]
One more word for describing what you lived and felt.
\end{grammarlens}

\subsection*{Your turn}
\begin{itemize}
\raggedright
  \item \emph{Say it:} \emph{kṣayaḥ}
  \item \emph{Say it:} \emph{kṣayaḥ}, once more
  \item \emph{Say it:} \emph{kṣayaḥ phupphusasya rogaḥ asti.} — Consumption is a disease of the lungs.
  \item \emph{Recall:} say the Sanskrit for a disease, then the Sanskrit for diabetes, then say \emph{kṣayaḥ} again
\end{itemize}

\begin{tcolorbox}[breakable,colback=teal!4,colframe=teal!35!black,title={Before you move on}]
What does \sk{क्षयः} mean? (\textquotedblleft{}A wasting disease\textquotedblright{}.) Say it once more.
\end{tcolorbox}
\section[\texorpdfstring{\sk{सञ्चरति} (sañcarati)}{सञ्चरति (sañcarati)}]{\sk{सञ्चरति} (sañcarati) — goes about; spreads}

\label{lesson:SA-C479-sancarati}



\begin{quote}
Before the new one: say the Sanskrit for diabetes, then the Sanskrit for a wasting disease.
\end{quote}

\subsection*{You'll want to know: \sk{सञ्चरति}}
\textbf{\sk{सञ्चरति}} — \emph{sañcarati} — \textquotedblleft{}goes about; spreads\textquotedblright{}.

\begin{quote}
\sk{ग्रीष्मे} \sk{ज्वरः} \sk{ग्रामे} \sk{सञ्चरति।}
\end{quote}

\begin{quote}
\emph{grīṣme jvaraḥ grāme sañcarati.}
\end{quote}

\begin{quote}
In summer fever goes about the village.
\end{quote}

Past (it …): \sk{ज्वरः} \sk{ग्रामे} \sk{सञ्चरितवान्}, \emph{jvaraḥ grāme sañcaritavān}.

\begin{cousinweb}
\textbf{\sk{सञ्चरति}} (\emph{sañcarati}) is built from \emph{sam}, \textquotedblleft{}together\textquotedblright{}, and the root √car, \textquotedblleft{}to move\textquotedblright{}, the root of \textbf{\sk{चरति}}; MW gives \textquotedblleft{}to go or wander about\textquotedblright{} (MW).
\end{cousinweb}

\begin{grammarlens}[title={What you've built}]
One more word for describing what you lived and felt.
\end{grammarlens}

\subsection*{Your turn}
\begin{itemize}
\raggedright
  \item \emph{Say it:} \emph{sañcarati}
  \item \emph{Say it:} \emph{sañcarati}, once more
  \item \emph{Say it:} \emph{grīṣme jvaraḥ grāme sañcarati.} — In summer fever goes about the village.
  \item \emph{Recall:} say the Sanskrit for diabetes, then the Sanskrit for a wasting disease, then say \emph{sañcarati} again
\end{itemize}

\begin{tcolorbox}[breakable,colback=teal!4,colframe=teal!35!black,title={Before you move on}]
What does \sk{सञ्चरति} mean? (\textquotedblleft{}Goes about; spreads\textquotedblright{}.) Say it once more.
\end{tcolorbox}
\section[\texorpdfstring{\sk{मारी} (mārī)}{मारी (mārī)}]{\sk{मारी} (mārī) — a plague}

\label{lesson:SA-C479-mari}



\begin{quote}
Before the new one: say the Sanskrit for a wasting disease, then the Sanskrit for goes about.
\end{quote}

\subsection*{You'll want to know: \sk{मारी}}
\textbf{\sk{मारी}} — \emph{mārī} — \textquotedblleft{}a plague\textquotedblright{}.

\begin{quote}
\sk{मार्याः} \sk{भयेन} \sk{जनाः} \sk{ग्रामं} \sk{त्यजन्ति।}
\end{quote}

\begin{quote}
\emph{māryāḥ bhayena janāḥ grāmaṃ tyajanti.}
\end{quote}

\begin{quote}
In fear of the plague people leave the village.
\end{quote}

A feminine noun.

\begin{cousinweb}
\textbf{\sk{मारी}} (\emph{mārī}) is built on \emph{māra}, \textquotedblleft{}killing\textquotedblright{}, from the root √mṛ, \textquotedblleft{}to die\textquotedblright{}, the root of \textbf{\sk{म्रियते}}; MW gives \textquotedblleft{}killing\textquotedblright{} and \textquotedblleft{}pestilence\textquotedblright{} (MW).
\end{cousinweb}

\begin{grammarlens}[title={What you've built}]
One more word for describing what you lived and felt.
\end{grammarlens}

\subsection*{Your turn}
\begin{itemize}
\raggedright
  \item \emph{Say it:} \emph{mārī}
  \item \emph{Say it:} \emph{mārī}, once more
  \item \emph{Say it:} \emph{māryāḥ bhayena janāḥ grāmaṃ tyajanti.} — In fear of the plague people leave the village.
  \item \emph{Recall:} say the Sanskrit for a wasting disease, then the Sanskrit for goes about, then say \emph{mārī} again
\end{itemize}

\begin{tcolorbox}[breakable,colback=teal!4,colframe=teal!35!black,title={Before you move on}]
What does \sk{मारी} mean? (\textquotedblleft{}A plague\textquotedblright{}.) Say it once more.
\end{tcolorbox}
\section[(dialogue)]{Looking back: from The Mouth to Diseases}

\label{lesson:SA-R479-review-472-479}



\begin{quote}
Say the Sanskrit for the palate and saliva.
\end{quote}

\subsection*{Your turn}
Say these aloud:

\begin{itemize}
\raggedright
  \item the palate, saliva, chews, spits, a taste
  \item a limb, the shin, the heel, a fist, bends a limb
  \item a vein, a sinew, a muscle, flesh, throbs
  \item sweats, yawns, sneezes, scratches, a hiccup
  \item strength, tall of stature, well-nourished, a figure, nourishes
  \item dizziness, vomits, grows weary, a burning feeling, runs a fever
  \item agony, coughs, dysentery, suffering, unbearable
  \item a disease, diabetes, a wasting disease, goes about, a plague
\end{itemize}

\begin{tcolorbox}[breakable,colback=teal!4,colframe=teal!35!black,title={Before you move on}]
The palate? (\textbf{\sk{तालु}}, \emph{tālu}.) Saliva? (\textbf{\sk{लाला}}, \emph{lālā}.) Chews? (\textbf{\sk{चर्वयति}}, \emph{carvayati}.) Spits? (\textbf{\sk{निष्ठीवति}}, \emph{niṣṭhīvati}.) A taste? (\textbf{\sk{रसः}}, \emph{rasaḥ}.) A limb? (\textbf{\sk{अंगम्}}, \emph{aṅgam}.) The shin? (\textbf{\sk{जंघा}}, \emph{jaṅghā}.) The heel? (\textbf{\sk{पार्ष्णिः}}, \emph{pārṣṇiḥ}.) A fist? (\textbf{\sk{मुष्टिः}}, \emph{muṣṭiḥ}.) Bends a limb? (\textbf{\sk{आकुञ्चयति}}, \emph{ākuñcayati}.) A vein? (\textbf{\sk{सिरा}}, \emph{sirā}.) A sinew? (\textbf{\sk{स्नायुः}}, \emph{snāyuḥ}.) A muscle? (\textbf{\sk{पेशी}}, \emph{peśī}.) Flesh? (\textbf{\sk{मांसम्}}, \emph{māṃsam}.) Throbs? (\textbf{\sk{स्पन्दते}}, \emph{spandate}.) Sweats? (\textbf{\sk{स्विद्यति}}, \emph{svidyati}.) Yawns? (\textbf{\sk{जृम्भते}}, \emph{jṛmbhate}.) Sneezes? (\textbf{\sk{क्षौति}}, \emph{kṣauti}.) Scratches? (\textbf{\sk{कण्डूयति}}, \emph{kaṇḍūyati}.) A hiccup? (\textbf{\sk{हिक्का}}, \emph{hikkā}.) Strength? (\textbf{\sk{शक्तिः}}, \emph{śaktiḥ}.) Tall of stature? (\textbf{\sk{प्रांशुः}}, \emph{prāṃśuḥ}.) Well-nourished? (\textbf{\sk{पुष्टः}}, \emph{puṣṭaḥ}.) A figure? (\textbf{\sk{आकृतिः}}, \emph{ākṛtiḥ}.) Nourishes? (\textbf{\sk{पोषयति}}, \emph{poṣayati}.) Dizziness? (\textbf{\sk{भ्रमः}}, \emph{bhramaḥ}.) Vomits? (\textbf{\sk{वमति}}, \emph{vamati}.) Grows weary? (\textbf{\sk{ग्लायति}}, \emph{glāyati}.) A burning feeling? (\textbf{\sk{दाहः}}, \emph{dāhaḥ}.) Runs a fever? (\textbf{\sk{ज्वरति}}, \emph{jvarati}.) Agony? (\textbf{\sk{वेदना}}, \emph{vedanā}.) Coughs? (\textbf{\sk{कासते}}, \emph{kāsate}.) Dysentery? (\textbf{\sk{अतिसारः}}, \emph{atisāraḥ}.) Suffering? (\textbf{\sk{क्लेशः}}, \emph{kleśaḥ}.) Unbearable? (\textbf{\sk{असह्यः}}, \emph{asahyaḥ}.) A disease? (\textbf{\sk{व्याधिः}}, \emph{vyādhiḥ}.) Diabetes? (\textbf{\sk{मधुमेहः}}, \emph{madhumehaḥ}.) A wasting disease? (\textbf{\sk{क्षयः}}, \emph{kṣayaḥ}.) Goes about? (\textbf{\sk{सञ्चरति}}, \emph{sañcarati}.) A plague? (\textbf{\sk{मारी}}, \emph{mārī}.)
\end{tcolorbox}
